An exact tensor-prime Euler transfer can avoid mixed primitive cycles by
making its atom graph acyclic.  We ask whether recurrence can instead be
retained in the base symbolic grammar while mixed lifts are removed by an
intrinsic trace.  Order the tensor atoms by entropy, put loops and both
nearest-neighbor arrows on the atom graph, label every directed cross edge
by a distinct positive generator of a free group, and assign the symmetric
endpoint weight
$a_n(s)=(p_n^{-s}+p_{n+1}^{-s})/2$.  The base graph is strongly connected
and aperiodic, but every mixed closed base word has a nonempty positive
cocycle word.  The canonical trace $\tau$ therefore kills it.  On the
semifinite algebra
$B(\ell^2(\At))\bar\otimes L(\mathbb F_\infty)$ we prove
\[
 (\Tr\otimes\tau)(T_s^r)=\sum_pp^{-rs},\qquad
 \det_\tau(I-zT_s)=\prod_p(1-zp^{-s})
\]
in the honest trace-class half-plane.  This defines SD-C10 and upgrades the
Euler ledger from an acyclic base to a recurrent one.  The mechanism is
positive-cone rather than specifically free: an abelian positive control
also works.  Its limitation is sharp.  The self-adjoint chiral double uses
$T_t^*$ and hence introduces inverse words.  Immediate backtracks
$gg^{-1}$ contribute the strictly positive mixed term
$4\sum_n|a_n(1/2+it)|^2$ to the quadratic trace.  This diverges, so neither
$\det_1$ nor $\det_2$ exists; $\det_3$ exists but subtracts precisely that
first recurrent term.  A two-atom formula proves that the fourth trace and
$\det_3$ still move with height.  Thus the analytic $\tau$-determinant sees
the Euler ledger but not the nonnormal geometry, while Fuglede--Kadison and
Brown data see magnitude geometry but do not supply a holomorphic Euler
divisor.  The stage reaches a recurrent noncommutative A2/A3 mechanism and
stops at a unified divisor.
